\documentclass[conference]{IEEEtran}
\usepackage{cite}
\usepackage{amsmath,amssymb,amsfonts}
\usepackage{algorithmic}
\usepackage{algorithm}
\usepackage{graphicx}
\usepackage{textcomp}
\usepackage{xcolor}
\usepackage{booktabs}
\usepackage{multirow}
\usepackage{array}
\usepackage{subcaption}
\usepackage{amsthm}
\usepackage{hyperref}
\usepackage{url}

% TikZ packages for diagrams
\usepackage{tikz}
\usetikzlibrary{positioning, shapes.geometric, arrows.meta, fit, backgrounds, decorations.pathreplacing, calc, shadows}

% Better table formatting
\usepackage{makecell}
\renewcommand\theadfont{\bfseries}
\renewcommand\theadalign{center}

% Setup hyperlinks
\hypersetup{
    colorlinks=true,
    linkcolor=blue,
    urlcolor=blue,
    citecolor=black,
}

% Theorem environments for privacy analysis
\newtheorem{theorem}{Theorem}
\newtheorem{definition}{Definition}
\newtheorem{lemma}{Lemma}

\begin{document}

\title{A Decentralized Privacy Preserving Smart Metering Protocol}

\author{
\IEEEauthorblockN{Chandrashekar Boddu}
\IEEEauthorblockA{Department of Computer Science\\
Lakehead University\\
Email: cboddu@lakeheadu.ca\\
\href{https://github.com/alexcs20/decentralized-smart-metering-protocol}{GitHub: github.com/alexcs20/smart-metering}}
\and
\IEEEauthorblockN{Dr. Mazhar Rathore}
\IEEEauthorblockA{Department of Computer Science\\
Lakehead University\\
Email: mazhar.rathore@lakeheadu.ca}
}

\maketitle

\begin{abstract}
A smart grid uses smart meters (SMs) to transmit high-frequency consumption data to utilities, enabling real-time monitoring and dynamic pricing. However, detailed consumption patterns risk exposing household activities, creating a critical privacy barrier. Privacy-preserving aggregation addresses this by reporting only aggregated data.

Existing aggregation approaches have limitations: Centralized frameworks risk a single point of compromise; in-network frameworks burden resource-constrained devices; and distributed frameworks often assume unrealistic adversary models or lack integrity verification and support for dynamic billing.

Our paper presents a decentralized protocol with three innovations. First, we use an Ethereum-style commit-reveal mechanism to generate unbiased, publicly verifiable randomness for selecting a temporary aggregation committee each round, eliminating the need for a permanent aggregator. Second, we combine Pedersen commitments with Shamir secret sharing to enable privacy-preserving aggregation with cryptographic integrity verification, detecting malicious share modifications by up to \(n-1\) compromised committee members. Third, our protocol integrates dynamic Time-of-Use (TOU) billing, supporting time-dependent rates while preserving individual reading privacy.

We implemented the protocol in C using GMP and OpenSSL, simulating 20 smart meters over 144 rounds (12 hours at 5-minute intervals) on Raspberry Pi 5 hardware. Consumption patterns reflected household types under Ontario's winter TOU rates. Results show lightweight operation: average message size is 254 bytes, and per-round computational overhead is 1.33 ms—favorable compared to 512–1024 bytes in homomorphic approaches. Committee selection was statistically fair (\(\chi^2 = 6.69 < 30.1\)), and the protocol scaled efficiently on embedded hardware.

Our protocol achieves an effective balance of privacy, integrity, efficiency, and scalability, making it suitable for real-world smart grid deployment.

\textbf{Code Availability:} \url{https://github.com/alexcs20/decentralized-smart-metering-protocol}
\end{abstract}

\begin{IEEEkeywords}
Smart grid, privacy preservation, committee selection, Pedersen commitments, Shamir secret sharing, Ethereum RANDAO
\end{IEEEkeywords}

\section{INTRODUCTION}

The electrical grid, originally designed for unidirectional power flow from centralized generators to passive consumers, is undergoing a fundamental transformation. The integration of digital communication technologies has given birth to the smart grid—a complex cyber-physical system where information flows bidirectionally alongside electricity. At the consumer level, this transformation is most visible through the deployment of smart meters (SMs), which have replaced traditional analog meters in millions of households worldwide. Unlike their predecessors, smart meters continuously record and transmit detailed consumption data, enabling utilities to monitor grid conditions in near real-time, implement dynamic pricing schemes, detect outages remotely, and optimize load balancing across the network \cite{tan2017survey}.

The promise of the smart grid extends beyond operational efficiency. By providing consumers with detailed feedback on their energy usage patterns, smart meters can encourage conservation and enable participation in demand response programs. Utilities gain unprecedented visibility into grid operations, allowing them to integrate variable renewable energy sources more effectively and defer costly infrastructure upgrades through better load management. These capabilities are essential for addressing the challenges of climate change and transitioning to a sustainable energy future, making the widespread deployment of smart metering infrastructure a critical priority for modern energy systems.

However, the same detailed consumption data that enables these benefits also creates significant privacy risks. Research has demonstrated that high-frequency meter readings can reveal far more than just total energy consumption. Through non-intrusive load monitoring techniques, analysts can disaggregate overall consumption into individual appliance signatures, revealing when residents cook, do laundry, watch television, or charge electric vehicles \cite{greveler2012multimedia}. More concerning still, studies have shown that it is possible to identify specific television programs or movies being watched based on the unique power consumption patterns of different content. In multi-occupancy buildings, occupancy patterns can be inferred, potentially enabling physical surveillance or burglary. These privacy implications represent a critical barrier to widespread smart grid adoption, necessitating privacy-preserving solutions that maintain grid functionality while protecting consumer data.

Privacy-preserving aggregation frameworks address these concerns by reporting only aggregated consumption to Electrical Service Providers (ESPs), thereby breaking the link between individual customers and their fine-grained consumption patterns. Rather than transmitting raw meter readings, these protocols enable the ESP to learn only aggregated consumption totals—sufficient for grid management and billing—while individual readings remain hidden. The challenge lies in designing protocols that achieve this privacy goal while satisfying the operational requirements of real-world smart grid deployments, including efficiency, scalability, and support for essential grid functionalities such as dynamic Time-of-Use (TOU) billing \cite{ben2022privacy, borges2014privacy, tan2024lppmm, zaredar2024lightweight}.

Existing approaches to privacy-preserving aggregation fall into three categories, each with distinct limitations. Centralized Aggregation Frameworks employ a single aggregator using homomorphic encryption but introduce a single point of compromise—if breached, all customer privacy is lost \cite{ben2022privacy}. In-network Aggregation Frameworks distribute aggregation across meter hierarchies but place computational burden on resource-constrained devices and remain vulnerable to attacks on parent nodes \cite{borges2014privacy}. Distributed Aggregation Frameworks leverage secret sharing to distribute trust across multiple aggregators, yet existing solutions either assume unrealistic semi-honest adversaries, lack integrity verification capabilities, or fail to support essential features such as dynamic TOU billing \cite{rottondi2013distributed}. More recent work by Tan et al. \cite{tan2024lppmm} proposed a lightweight privacy-preserving scheme for multi-dimensional and multi-subset data aggregation, while Zaredar and Amini \cite{zaredar2024lightweight} developed a privacy-preserving billing protocol with dynamic tariff policy adjustment. Additionally, research on privacy-preserving data provenance \cite{parameswarath2024provenance}, real-time privacy-preserving data release \cite{shateri2024realtime}, hybrid obfuscation schemes \cite{tran2024hybrid}, and streaming data privacy against appliance status leakage \cite{hong2018privacy} have explored various aspects of smart meter privacy. Furthermore, Yao et al. \cite{yao2024theft} investigated the intersection of energy theft detection and privacy preservation, while Mustafa et al. \cite{mustafa2023secure} proposed a secure and privacy-preserving protocol for operational data collection. Comprehensive systematizations of knowledge in this domain \cite{behrouz2024privacy} have provided valuable taxonomies and frameworks for understanding the landscape of smart meter privacy techniques. Complementary to privacy concerns, decentralized coordination of household energy flexibility \cite{wang2024decentralized} has emerged as an important paradigm for virtual power plants, highlighting the growing need for decentralized architectures. However, these approaches still rely on centralized or semi-centralized architectures. These limitations highlight the need for a decentralized approach that provides strong privacy guarantees, cryptographic integrity verification, and practical functionality under realistic threat models.

Our paper presents a novel decentralized protocol that addresses these limitations through three key innovations. First, we employ an Ethereum-style commit-reveal mechanism to generate unbiased, publicly verifiable randomness for selecting a temporary aggregation committee each round, completely eliminating the need for a permanent aggregator while ensuring unpredictable, fair, and verifiable committee selection. Second, we combine Pedersen commitments \cite{pedersen1991non} with Shamir secret sharing \cite{shamir1979how} to enable privacy-preserving aggregation with cryptographic integrity verification, detecting any malicious modification of shares by up to \(n-1\) compromised committee members—even under a dishonest majority. Third, our protocol seamlessly integrates dynamic TOU billing, enabling accurate customer billing based on time-dependent rates while preserving individual reading privacy.

To validate real-world deployment viability, we implemented the protocol in C using GMP for large integer arithmetic and OpenSSL for cryptographic primitives, conducting extensive simulations with 20 smart meters over 144 rounds (12 hours of operation with 5-minute intervals) on Raspberry Pi 5 hardware. Results demonstrate lightweight operation with average message sizes of 254 bytes and per-round computational overhead of just 1.33 ms—significantly smaller than the 512–1024 bytes typical of homomorphic approaches. The committee selection mechanism proves statistically fair, and the protocol scales gracefully on resource-constrained embedded hardware. Collectively, our contributions demonstrate that the protocol achieves an effective balance of privacy, integrity, efficiency, and scalability suitable for real-world smart grid deployment.

\section{RELATED WORK}

Privacy-preserving aggregation for smart metering has been extensively studied, with frameworks generally classified into three categories based on architectural choices.

\subsection{Centralized Aggregation Frameworks (CAF)}

Centralized frameworks employ a single aggregator that collects encrypted readings from all smart meters and reports only aggregated totals to the Electrical Service Provider (ESP). Efthymiou and Kalogridis introduced anonymization using pseudonyms, where meters send readings under temporary identifiers to separate identity from consumption data. Garcia and Jacobs first proposed using Paillier homomorphic encryption for smart meter aggregation, enabling the aggregator to compute encrypted sums without accessing individual readings. Ben Romdhane et al. \cite{ben2022privacy} extended this approach by combining Paillier encryption with ECDHE key establishment, adding ephemeral session keys for each reporting round. Lu et al. developed EPPA, which uses superincreasing sequences to aggregate multi-dimensional data in a single encrypted message. More recently, Tan et al. \cite{tan2024lppmm} proposed LPPMM-DA, a lightweight privacy-preserving scheme for multi-dimensional and multi-subset data aggregation that improves efficiency while maintaining strong privacy guarantees. Additionally, Zaredar and Amini \cite{zaredar2024lightweight} developed a lightweight privacy-preserving billing protocol that supports dynamic tariff policy adjustment, enabling flexible pricing models. However, these centralized approaches remain vulnerable to a single point of compromise.

\subsection{In-network Aggregation Frameworks (IAF)}

In-network frameworks distribute aggregation across smart meters organized in a hierarchical structure, eliminating the need for a dedicated aggregator. Borges et al. \cite{borges2014privacy} developed a tree-based protocol where meters encrypt readings using Paillier, and parent nodes aggregate ciphertexts from their children before forwarding results upward. The root node sends the final aggregated ciphertext to the ESP for decryption. Their scheme also incorporates commitments to enable verification. Erkin and Tsudik proposed using additive secret sharing in a hierarchical structure, where meters split readings into shares distributed along the tree, with intermediate nodes combining shares before forwarding. Li and Luo introduced integrity verification using Merkle hash trees, enabling detection of data modification during aggregation.

\subsection{Distributed and Privacy-Preserving Frameworks}

Distributed frameworks leverage secret sharing to distribute trust across multiple Dedicated Aggregators (DAs), ensuring that no single entity can reconstruct individual readings. Rottondi et al. \cite{rottondi2013distributed} pioneered the use of Shamir secret sharing for smart metering, where each meter splits its reading into shares distributed among \(n\) DAs using a \((k,n)\) threshold scheme. The DAs independently aggregate received shares, and any \(k\) aggregated shares can reconstruct the total consumption. Garcia et al. incorporated differential privacy by adding calibrated noise to shares, providing formal privacy guarantees against inference attacks. Wagh et al. combined Pedersen commitments \cite{pedersen1991non} with Shamir secret sharing \cite{shamir1979how} to enable integrity verification, allowing detection of share modification through commitment aggregation. Guan et al. explored blockchain-based coordination, where meters submit encrypted readings to a blockchain and a committee performs aggregation through smart contracts, providing transparency and immutability.

Beyond aggregation, several complementary privacy-preserving techniques have been proposed for smart meter communications. Parameswarath and Sikdar \cite{parameswarath2024provenance} introduced a privacy-preserving data provenance scheme that enables verification of data origin without compromising consumer privacy, which is complementary to our integrity verification approach. Shateri et al. \cite{shateri2024realtime} developed a real-time privacy-preserving data release mechanism that balances utility and privacy in streaming meter data. Tran et al. \cite{tran2024hybrid} proposed a hybrid obfuscation scheme that operates without a trusted third party, combining multiple privacy techniques to protect against inference attacks. Hong et al. \cite{hong2018privacy} analyzed information leakage of appliance status from smart meter streaming data and proposed countermeasures to mitigate such leakage. Yao et al. \cite{yao2024theft} investigated the challenging problem of energy theft detection while preserving consumer privacy, demonstrating that privacy and utility can be balanced effectively. Mustafa et al. \cite{mustafa2023secure} designed a secure and privacy-preserving protocol specifically for smart metering operational data collection, addressing both confidentiality and integrity requirements.

A comprehensive systematization of knowledge by Behrouz et al. \cite{behrouz2024privacy} provides a structured taxonomy of privacy-preserving methods for smart meter data, categorizing approaches by their underlying techniques, threat models, and privacy guarantees. This work serves as an important reference for understanding the landscape of smart meter privacy research. In parallel with privacy concerns, the decentralization of energy management has gained significant attention, with Wang et al. \cite{wang2024decentralized} proposing a decentralized coordination framework for household energy flexibility in virtual power plants, highlighting the growing need for decentralized architectures in smart grid systems. Our work builds upon distributed aggregation principles while introducing Ethereum-style committee selection that rotates the aggregation responsibility each round, eliminating the need for static aggregators entirely.

\section{CRYPTOGRAPHIC PRELIMINARIES}

This section introduces the fundamental cryptographic primitives, notation, and concepts that form the foundation of our protocol.

\subsection{Notation and Terminology}

Table~\ref{tab:notation} summarizes the notation used throughout this paper.

\begin{table}[h]
\centering
\caption{Summary of Notation}
\label{tab:notation}
\begin{tabular}{|c|l|}
\hline
\textbf{Symbol} & \textbf{Description} \\
\hline
$N$ & Total number of smart meters \\
$SM_i$ & Smart meter with identifier $i$, where $i \in \{1,2,\ldots,N\}$ \\
$c$ & Committee size (number of meters selected per round) \\
$k$ & Reconstruction threshold (minimum shares needed) \\
$r$ & Round number \\
$R_{i,r}$ & Energy reading of $SM_i$ in round $r$ \\
$\rho_i$ & Random number generated by $SM_i$ for commit-reveal \\
$H(\cdot)$ & Cryptographic hash function (SHA-256) \\
$\parallel$ & Concatenation operator \\
$seed_r$ & Random seed for round $r$ \\
$C_{i,r}$ & Pedersen commitment of $SM_i$ for round $r$ \\
$g, h$ & Public generators for Pedersen commitments \\
$p$ & Large prime modulus \\
$r_{i,r}$ & Random blinding factor for commitment \\
$F_{i,r}(x)$ & Secret sharing polynomial of $SM_i$ for round $r$ \\
$a_{i,d,r}$ & Coefficient of degree $d$ in polynomial $F_{i,r}(x)$ \\
$s_{i\to j,r}$ & Share sent from $SM_i$ to committee member $j$ in round $r$ \\
$S_{j,r}$ & Aggregated share held by committee member $j$ \\
$SCR_r$ & Aggregated spatial commitment register for round $r$ \\
$\delta_j(x)$ & Lagrange basis polynomial for committee member $j$ \\
$M_r$ & Total energy consumption in round $r$ (spatial aggregate) \\
$\beta_{d,r}$ & Coefficient of reconstructed polynomial for round $r$ \\
$M_{i,T}$ & Temporal aggregate for $SM_i$ over period $T$ \\
$p_r$ & TOU electricity price for round $r$ \\
$Bill_i$ & Final bill for customer $i$ \\
$\chi^2$ & Chi-square statistic for fairness testing \\
\hline
\end{tabular}
\end{table}

\subsection{Pedersen Commitments}

Pedersen commitments \cite{pedersen1991non} provide a cryptographic primitive that allows a party to commit to a secret value while keeping it hidden, with the ability to later reveal the committed value. The scheme is information-theoretically hiding and computationally binding under the discrete logarithm assumption. For a cyclic group $\mathbb{G}$ of prime order $p$ with generators $g$ and $h$ where the discrete logarithm of $h$ base $g$ is unknown to all parties, a commitment to a message $m \in \mathbb{Z}_p$ is computed as

\begin{equation}
C = g^{m} \cdot h^{r} \pmod{p}
\end{equation}

where $r \in \mathbb{Z}_p$ is a randomly chosen blinding factor. The commitment $C$ reveals no information about $m$ because for any potential message $m'$, there exists a corresponding $r'$ that would produce the same commitment. Crucially for our protocol, Pedersen commitments are homomorphic under multiplication

\begin{equation}
C(m_1, r_1) \cdot C(m_2, r_2) = g^{m_1+m_2} \cdot h^{r_1+r_2} = C(m_1 + m_2, r_1 + r_2)
\end{equation}

This property enables the aggregation of commitments without decrypting individual values, forming the basis of our integrity verification mechanism.

\subsection{Shamir Secret Sharing}

Shamir's \((k,n)\) threshold secret sharing scheme \cite{shamir1979how} enables a dealer to split a secret $s$ into $n$ shares such that any $k$ shares can reconstruct the secret, but any $k-1$ shares reveal no information about it. The scheme operates on a finite field $\mathbb{F}_p$ of prime order $p$. The dealer constructs a random polynomial of degree $k-1$

\begin{equation}
F(x) = s + a_1x + a_2x^2 + \ldots + a_{k-1}x^{k-1} \pmod{p}
\end{equation}

where $a_1, a_2, \ldots, a_{k-1}$ are uniformly random coefficients. Shares are computed as $F(i)$ for $i = 1, 2, \ldots, n$ and distributed to the $n$ participants. Reconstruction uses Lagrange interpolation to recover the polynomial $F(x)$ from any $k$ shares $(x_j, y_j)$

\begin{equation}
F(0) = \sum_{j=1}^{k} y_j \cdot \delta_j(0)
\end{equation}

where the Lagrange basis polynomials are defined as

\begin{equation}
\delta_j(x) = \prod_{\substack{1 \leq m \leq k \\ m \neq j}} \frac{x - x_m}{x_j - x_m}
\end{equation}

In our protocol, each smart meter acts as a dealer, splitting its consumption reading $R_{i,r}$ into $c$ shares distributed to committee members. Notably, we set the first coefficient $a_1$ equal to the blinding factor $r_{i,r}$ used in the Pedersen commitment, creating a cryptographic link that enables integrity verification.

\subsection{Ethereum-Style Committee Selection}

Randomness generation and unbiased committee selection are critical components of decentralized protocols. Ethereum's RANDAO mechanism provides an elegant solution for generating publicly verifiable randomness without trusting any single party. The protocol operates by having each participant commit to a random value $\rho_i$ by publishing $H(r \parallel i \parallel \rho_i)$. After all commitments are received, participants reveal their $\rho_i$ values. The final seed is computed as

\begin{equation}
\mathit{seed}_r = H(r \parallel \rho_1 \parallel \rho_2 \parallel \ldots \parallel \rho_q)
\end{equation}

This approach ensures that the randomness is unpredictable because a participant cannot change their $\rho_i$ after seeing others' commitments. The final seed is verifiable by all participants, as anyone can recompute the hash chain. Our protocol adapts this mechanism for committee selection by using the seed to generate deterministic scores

\begin{equation}
\mathit{score}_{i,j} = H(\mathit{seed}_r \parallel j)
\end{equation}

The $c$ participants with the smallest scores form the aggregation committee for round $r$.

\subsection{Threat Model and Security Assumptions}

We consider a malicious adversary capable of compromising up to $n-1$ committee members in any round. The adversary may modify or drop shares sent to committee members, alter aggregated shares during reconstruction, attempt to reconstruct individual readings by colluding with other committee members, or influence randomness generation by selectively revealing $\rho_i$ values. We assume that the discrete logarithm problem is computationally infeasible in the chosen group, smart meters are trusted and correctly execute the protocol, the ESP is honest-but-curious (following protocol specifications but attempting to learn additional information from observed data), and communication channels are authenticated to ensure messages cannot be spoofed. The protocol's security guarantees hold under these assumptions, providing privacy and integrity even when a dishonest majority of committee members are compromised.

\section{PROPOSED PROTOCOL}

\subsection{System Model}

The system consists of $N$ smart meters $\{SM_1, SM_2, \ldots, SM_N\}$ and an Electrical Service Provider (ESP). Time is divided into discrete rounds $r \in \{1,2,\ldots,R\}$ with each round representing a fixed time interval (e.g., 5 minutes). In each round, a committee $C_r$ of size $c$ is temporarily selected to perform aggregation on behalf of all meters. The protocol employs a $(k,n)$ threshold secret sharing scheme where $c = 10$ and the reconstruction threshold is $k = 6$, meaning any 6 of the 10 committee members can reconstruct the aggregated total. The protocol operates under the key principles of having no permanent aggregator to eliminate any single point of trust or compromise, preserving privacy such that individual meter readings are never revealed to any party including committee members and the ESP, providing integrity verification so that any malicious modification of shares during aggregation is cryptographically detectable, and supporting dynamic billing natively without additional privacy compromises.

\subsection{Protocol Overview}

Our protocol operates in eight distinct phases that collectively achieve privacy-preserving aggregation with integrity verification. As illustrated in Figure~\ref{fig:protocol_flow}, these phases can be grouped into three logical stages: randomness generation and committee selection, data submission and secure sharing, and aggregation, reconstruction, and reporting.

\begin{figure}[ht]
\centering
\includegraphics[width=0.9\columnwidth]{protocol_flow.png}
\caption{Smart meter commitment and aggregation flow: Phases 1-2 generate public randomness and select committee; Phases 3-4 handle Pedersen commitments and Shamir share distribution; Phases 5-6 perform local aggregation; Phases 7-8 reconstruct totals, verify integrity, and compute TOU bills.}
\label{fig:protocol_flow}
\end{figure}

In the first stage comprising Phases 1 and 2, the protocol begins with a distributed randomness generation process where all meters contribute to an unbiased, publicly verifiable random seed. This process uses an Ethereum-style commit-reveal mechanism that prevents any single meter from influencing or predicting the final seed. Once the seed is established, each meter independently computes deterministic scores for all meters using the seed as input. The $c$ meters with the smallest scores are selected as the aggregation committee for the current round. This approach ensures that no meter can know the committee composition before the seed is revealed, every meter has an equal probability of being selected in each round, any participant can independently recompute and verify the committee selection, and a completely new committee is formed each round preventing any permanent concentration of power.

In the second stage comprising Phases 3 and 4, once the committee is established, each meter commits to its consumption reading using Pedersen commitments \cite{pedersen1991non}. The Pedersen commitment scheme provides information-theoretic hiding, meaning the commitment reveals nothing about the actual reading, while being computationally binding—the meter cannot later claim a different reading. Each meter also generates a random blinding factor that serves dual purposes: it masks the commitment and serves as the first coefficient in the secret sharing polynomial. Each meter then splits its reading using Shamir's $(k,c)$ threshold secret sharing scheme \cite{shamir1979how}. The meter constructs a random polynomial of degree $k-1$ where the constant term is the actual reading and the first coefficient is the blinding factor. The meter computes $c$ shares by evaluating this polynomial at distinct points and sends one share to each committee member. This distribution ensures that no single committee member receives enough information to reconstruct any individual reading, any $k$ committee members can later reconstruct the total consumption through Lagrange interpolation, and the blinding factor embedded in the polynomial enables cryptographic integrity verification.

In the third stage comprising Phases 5 through 8, committee members locally aggregate the shares they receive from all meters. Due to the additive homomorphism of Shamir secret sharing, the sum of the shares received by a committee member equals the share of the total consumption that member would hold if all meters had directly shared their readings. The committee also aggregates all Pedersen commitments by multiplying them, leveraging the homomorphic property to obtain a single commitment representing the total consumption. The committee members broadcast their aggregated shares. Once any $k$ members have revealed their shares, the ESP can reconstruct the total consumption using Lagrange interpolation. The ESP then verifies integrity by checking whether the reconstructed total matches the aggregated commitment using the verification equation. This verification detects any malicious modification of shares by up to $n-1$ committee members. Finally, the ESP performs temporal aggregation across multiple rounds for billing purposes. The protocol supports dynamic Time-of-Use billing by applying the appropriate rate to each round's consumption before aggregation, similar to the dynamic tariff adjustment mechanisms explored in \cite{zaredar2024lightweight}. The ESP learns only the final bill for each customer, never the individual per-round readings.

\subsection{Phase Details}

\subsubsection{Phase 1: Commit-Reveal (Randomness Generation)}

Each smart meter $SM_i$ generates a random number $\rho_i$ from a uniform distribution over a large domain. It then broadcasts a commitment to this random number

\begin{equation}
\mathit{commit}_i = H(r \parallel i \parallel \rho_i)
\end{equation}

where $H(\cdot)$ is a cryptographic hash function (SHA-256). This commitment hides $\rho_i$ while binding $SM_i$ to its chosen value. After all commitments are received, meters reveal their $\rho_i$ values. Each reveal is verified by recomputing the hash and comparing with the previously broadcast commitment. The round seed is computed as

\begin{equation}
\mathit{seed}_r = H(r \parallel \rho_1 \parallel \rho_2 \parallel \ldots \parallel \rho_q)
\end{equation}

where $q$ is the number of valid reveals. This seed is unpredictable until the reveal phase, as a malicious meter cannot change its $\rho_i$ after seeing others' commitments.

\subsubsection{Phase 2: Ethereum-Style Committee Selection}

Using the seed, each meter independently computes a deterministic score for every meter

\begin{equation}
\mathit{score}_{i,j} = H(\mathit{seed}_r \parallel j)
\end{equation}

The $c$ meters with the smallest scores form the committee $C_r$ for round $r$. This selection mechanism guarantees randomness because the committee composition is a deterministic function of the random seed, fairness because each meter has equal probability of selection as hash outputs are uniformly distributed, and verifiability because any participant can recompute the scores and verify the committee composition.

\subsubsection{Phase 3: Pedersen Commitments}

Each meter $SM_i$ generates a Pedersen commitment \cite{pedersen1991non} for its reading $R_{i,r}$

\begin{equation}
C_{i,r} = g^{R_{i,r}} \cdot h^{r_{i,r}} \pmod{p}
\end{equation}

where $r_{i,r}$ is a random blinding factor chosen uniformly from $\mathbb{Z}_p$. The generators $g$ and $h$ are public system parameters with the discrete logarithm of $h$ base $g$ unknown. The commitment is broadcast to all participants and will be used later for integrity verification.

\subsubsection{Phase 4: Share Generation and Distribution}

Each meter constructs a $(k-1)$-degree polynomial over $\mathbb{F}_p$

\begin{equation}
F_{i,r}(x) = R_{i,r} + a_{i,1}x + a_{i,2}x^2 + \ldots + a_{i,k-1}x^{k-1}
\end{equation}

where the constant term is the reading $R_{i,r}$ and the first coefficient $a_{i,1} = r_{i,r}$—the same blinding factor used in the commitment. The remaining coefficients $a_{i,2}, \ldots, a_{i,k-1}$ are chosen uniformly at random. Shares are computed for each committee member $j \in C_r$ as

\begin{equation}
s_{i \to j,r} = F_{i,r}(j)
\end{equation}

Each share is encrypted using the committee member's public key and transmitted securely.

\subsubsection{Phase 5: Spatial Aggregation at Committee}

Upon receiving shares from all meters, each committee member $j$ computes the sum of all received shares

\begin{equation}
S_{j,r} = \sum_{i=1}^{N} s_{i \to j,r}
\end{equation}

Due to the linearity of polynomial evaluation, $S_{j,r}$ equals the evaluation at point $j$ of the sum polynomial $\sum_{i=1}^{N} F_{i,r}(x)$. The committee member also aggregates all received commitments

\begin{equation}
\mathit{SCR}_r = \prod_{i=1}^{N} C_{i,r} \pmod{p}
\end{equation}

By the homomorphic property of Pedersen commitments \cite{pedersen1991non}, this equals $C(\sum R_{i,r}, \sum r_{i,r})$.

\subsubsection{Phase 6: Reconstruction of Spatial Aggregate}

Committee members broadcast their aggregated shares $S_{j,r}$. Once any $k$ shares are available, the ESP reconstructs the total consumption using Lagrange interpolation

\begin{equation}
M_r = \sum_{j \in C_r} S_{j,r} \cdot \delta_j(0) \pmod{p}
\end{equation}

where $\delta_j(x)$ are the Lagrange basis polynomials. The ESP also reconstructs the sum of blinding factors $\beta_{1,r}$ using the same interpolation, as it is the first coefficient of the reconstructed sum polynomial. The ESP verifies integrity by checking

\begin{equation}
g^{M_r} \cdot h^{\beta_{1,r}} \stackrel{?}{=} \mathit{SCR}_r
\end{equation}

If this equality holds, it proves that all meters correctly contributed their shares and that no shares were modified during aggregation. Any malicious modification by up to $n-1$ committee members will cause this verification to fail.

\subsubsection{Phase 7: Temporal Aggregation}

For billing period $T$ (e.g., one month), the ESP computes each customer's temporal aggregate

\begin{equation}
M_{i,T} = \sum_{r \in T} R_{i,r}
\end{equation}

For Time-of-Use billing, the ESP applies time-dependent prices $p_r$

\begin{equation}
\mathit{Bill}_i = \sum_{r \in T} R_{i,r} \cdot p_r
\end{equation}

The ESP learns only the final bill amount, not the individual per-round readings. This dynamic billing approach aligns with the flexible tariff adjustment mechanisms studied in \cite{zaredar2024lightweight}.

\subsubsection{Phase 8: Reporting to ESP}

The committee reports only the final spatial aggregate $M_r$ and the verification parameters to the ESP. Individual readings, shares, and commitments are never revealed to any party. The ESP stores only the aggregated totals and uses them for billing and grid management.

\subsection{Security Properties}

The protocol provides several important security properties. Regarding privacy, individual readings are never revealed to any party. Committee members only see shares, which individually provide no information about the original reading due to the properties of Shamir secret sharing \cite{shamir1979how}. The ESP sees only aggregated totals. For integrity, any malicious modification of shares by committee members is detectable through the commitment verification equation, and the protocol can identify up to $n-1$ malicious committee members. Regarding fairness, committee selection is unbiased and verifiable, with each meter having equal probability of selection in each round. For collusion resistance, at least $k$ committee members must collude to reconstruct individual readings, providing strong protection against dishonest majority attacks.

\section{EXPERIMENTAL RESULTS}

\subsection{Experimental Setup}

We implemented our protocol in C using the GMP library for large integers and OpenSSL for cryptographic primitives. The complete implementation is available online (see Code Availability section).

Simulations were run for 144 rounds (12 hours) with 20 smart meters on a Raspberry Pi 5 (4GB RAM, 2.4GHz quad-core ARM Cortex-A76). Each meter generated realistic consumption patterns based on household type. Time-of-Use rates followed Ontario's winter tariff structure (off-peak: 9.8\textcent/kWh, mid-peak: 15.7\textcent/kWh, on-peak: 20.3\textcent/kWh).

\begin{table}[h]
\centering
\caption{Simulation Parameters}
\label{tab:params}
\begin{tabular}{|l|c|}
\hline
\textbf{Parameter} & \textbf{Value} \\
\hline
Platform & Raspberry Pi 5 \\
Number of Meters ($N$) & 20 \\
Committee Size ($c$) & 10 \\
Threshold ($k$) & 6 \\
Rounds & 144 \\
Prime Size ($p$) & 512 bits \\
\hline
\end{tabular}
\end{table}

\subsection{Computational Overhead}

Table II presents the average computational overhead per round for each protocol phase measured on the Raspberry Pi 5. Phase 1 (Commit-Reveal) requires 322 microseconds representing 24.2\% of the total overhead, Phase 2 (Committee Selection) requires 167 microseconds representing 12.6\%, Phase 3 (Pedersen Commitments) requires 185 microseconds representing 13.9\%, Phase 4 (Share Generation) requires 297 microseconds representing 22.4\%, Phase 5 (Spatial Aggregation) requires 53 microseconds representing 4.0\%, Phase 6 (Reconstruction) requires 153 microseconds representing 11.5\%, and Phase 8 (Billing) requires 152 microseconds representing 11.4\%. The total per-round overhead is 1,329 microseconds.

\begin{table}[h]
\centering
\caption{Average Phase Timing (microseconds) - Raspberry Pi 5}
\label{tab:timing}
\begin{tabular}{|l|c|c|}
\hline
\textbf{Phase} & \textbf{Time ($\mu$s)} & \textbf{\% of Total} \\
\hline
Phase 1 (Commit-Reveal) & 322 & 24.2\% \\
Phase 2 (Committee Selection) & 167 & 12.6\% \\
Phase 3 (Pedersen Commitments) & 185 & 13.9\% \\
Phase 4 (Share Generation) & 297 & 22.4\% \\
Phase 5 (Spatial Aggregation) & 53 & 4.0\% \\
Phase 6 (Reconstruction) & 153 & 11.5\% \\
Phase 8 (Billing) & 152 & 11.4\% \\
\hline
\textbf{Total per Round} & \textbf{1,329} & \textbf{100\%} \\
\hline
\end{tabular}
\end{table}

\subsection{Communication Overhead}

The communication metrics from our Raspberry Pi simulation show that total messages sent were 23,184, total data transferred was 5.76 MB, average message size was 254 bytes, minimum message size was 128 bytes, and maximum message size was 512 bytes. The average message size of 254 bytes is significantly smaller than homomorphic encryption-based approaches which typically require 512 to 1024 bytes, making our protocol suitable for resource-constrained smart meters.

\begin{table}[h]
\centering
\caption{Communication Metrics - Raspberry Pi Simulation}
\label{tab:comm}
\begin{tabular}{|l|c|}
\hline
\textbf{Metric} & \textbf{Value} \\
\hline
Total Messages Sent & 23,184 \\
Total Data Transferred & 5.76 MB \\
Average Message Size & 254 bytes \\
Min Message Size & 128 bytes \\
Max Message Size & 512 bytes \\
\hline
\end{tabular}
\end{table}

\subsection{Performance Scaling Analysis}

Performance scaling projections for larger networks show that for 20 meters the total per-round overhead is 1.33 ms, for 100 meters it is 4.44 ms, and for 500 meters it is 20.44 ms. This demonstrates that our protocol scales efficiently due to distributed computation across committee members, parallel share processing, and minimal communication rounds requiring only one broadcast per phase.

\begin{table}[h]
\centering
\caption{Performance Scaling Projections for Larger Networks}
\label{tab:scaling}
\begin{tabular}{|l|c|c|c|}
\hline
\textbf{Component} & \textbf{20 Meters} & \textbf{100 Meters} & \textbf{500 Meters} \\
\hline
Commit-Reveal & 322 $\mu$s & 1.61 ms & 8.05 ms \\
Committee Selection & 167 $\mu$s & 167 $\mu$s & 167 $\mu$s \\
Pedersen Commitments & 185 $\mu$s & 925 $\mu$s & 4.63 ms \\
Share Generation & 297 $\mu$s & 1.49 ms & 7.43 ms \\
Spatial Aggregation & 53 $\mu$s & 53 $\mu$s & 53 $\mu$s \\
Reconstruction & 53 $\mu$s & 53 $\mu$s & 53 $\mu$s \\
Billing & 52 $\mu$s & 52 $\mu$s & 52 $\mu$s \\
\hline
\textbf{Total per Round} & \textbf{1.33 ms} & \textbf{4.44 ms} & \textbf{20.44 ms} \\
\hline
\end{tabular}
\end{table}

\subsection{Committee Selection Fairness Analysis}

The committee selection statistics show that committee changes occur every round (286 total changes), average formation time is 134 microseconds, selection range per meter is between 137 and 149 times, theoretical expected selections per meter is 72, and the maximum versus minimum difference is 8.7\%. The committee selection mechanism achieves statistical fairness, as evidenced by the chi-square test

\begin{equation}
\chi^2 = \sum_{i=1}^{N} \frac{(\mathit{observed}_i - \mathit{expected}_i)^2}{\mathit{expected}_i} = 6.69
\end{equation}

With $N-1 = 19$ degrees of freedom, the critical value at $\alpha = 0.05$ is 30.1. Our $\chi^2 = 6.69$ confirms unbiased selection.

\begin{table}[h]
\centering
\caption{Committee Selection Statistics}
\label{tab:committee}
\begin{tabular}{|l|c|}
\hline
\textbf{Metric} & \textbf{Value} \\
\hline
Committee Changes & 286 (every round) \\
Average Formation Time & 134 $\mu$s \\
Selection Range (per meter) & 137-149 times \\
Theoretical Expected & 72 selections per meter \\
Max vs Min Difference & 8.7\% \\
\hline
\end{tabular}
\end{table}

\section{SECURITY ANALYSIS}

\subsection{Privacy Guarantees}

The protocol provides strong privacy guarantees through multiple mechanisms. Individual readings are never revealed to any party because the ESP only receives aggregated values and committee members only see shares, which individually reveal nothing about the original reading due to the properties of Shamir secret sharing \cite{shamir1979how}. Unlinkability is achieved through rotating committees that ensure no single entity can link readings across time, as each round uses a completely new committee. Perfect hiding is provided by Pedersen commitments \cite{pedersen1991non}, which offer information-theoretic hiding of meter readings, meaning even an adversary with unlimited computational power cannot extract the reading from the commitment.

\subsection{Integrity Verification}

The protocol detects any modification to shares through the commitment verification equation. If shares are modified before aggregation, the reconstructed polynomial will not match the aggregated commitment, triggering detection. The verification works even when up to $n-1$ committee members are compromised, as the honest committee members' shares will reveal the discrepancy.

\subsection{Collusion Resistance}

With threshold $k = 6$, at least 6 of 10 committee members must collude to reconstruct individual readings. This provides strong protection against dishonest majority attacks, as the adversary would need to compromise a significant portion of the committee to breach privacy.

\section{PRIVACY ANALYSIS}

This section provides a formal analysis of the privacy guarantees offered by our protocol, examining resistance against various attack vectors and comparing with existing privacy-preserving approaches.

\subsection{Privacy Threat Model}

We consider an adversary with the following capabilities: (1) compromise of up to \(k-1\) committee members in any round, (2) passive eavesdropping on all communication channels, (3) ability to influence randomness generation through selective reveal strategies, and (4) possession of auxiliary information including household characteristics and historical consumption patterns. The adversary's objective is to infer individual meter readings \(R_{i,r}\) for specific households or time periods.

The protocol provides privacy guarantees under the following assumptions: the discrete logarithm problem is computationally infeasible in the chosen group \(\mathbb{G}\) of order \(p\), the hash function \(H(\cdot)\) behaves as a random oracle, and the underlying communication channels are authenticated to prevent message spoofing.

\subsection{Information-Theoretic Privacy Analysis}

\subsubsection{Share Indistinguishability}

Shamir's \((k,c)\) secret sharing scheme \cite{shamir1979how} provides perfect secrecy: any set of fewer than \(k\) shares reveals no information about the secret. For our protocol, each committee member receives exactly one share from each smart meter. The share \(s_{i \to j,r}\) held by committee member \(j\) is a point on the random polynomial \(F_{i,r}(x)\) of degree \(k-1\).

\begin{theorem}
For any adversary that compromises fewer than \(k\) committee members, the joint distribution of compromised shares is independent of the meter reading \(R_{i,r}\).
\end{theorem}

\begin{proof}
The polynomial \(F_{i,r}(x) = R_{i,r} + a_{i,1}x + \ldots + a_{i,k-1}x^{k-1}\) has coefficients chosen uniformly at random from \(\mathbb{F}_p\). For any set of \(t < k\) evaluation points \(\{j_1, \ldots, j_t\}\), the mapping from coefficients \((R_{i,r}, a_{i,1}, \ldots, a_{i,k-1})\) to shares \((F(j_1), \ldots, F(j_t))\) is a linear transformation with a \((k-t)\)-dimensional kernel. Each possible secret value \(R'\) corresponds to a distinct affine subspace of the same size, giving equal probability for all secrets. Thus, \(\Pr[R_{i,r} = r \mid \text{compromised shares}] = \Pr[R_{i,r} = r] = 1/p\).
\end{proof}

With threshold \(k = 6\), an adversary compromising up to 5 committee members learns nothing about any individual meter reading. Even with auxiliary information about household characteristics, the adversary's posterior distribution over possible readings remains uniform.

\subsubsection{Pedersen Commitment Hiding}

Pedersen commitments \cite{pedersen1991non} provide information-theoretic hiding: for any commitment \(C = g^m h^r\), and any candidate message \(m'\), there exists a unique \(r'\) such that \(g^{m'} h^{r'} = C\). This holds because \(h\) is a generator of the group, and the discrete logarithm between \(g\) and \(h\) is unknown.

\begin{equation}
\forall m, m' \in \mathbb{F}_p, \exists! r' : g^{m} h^{r} = g^{m'} h^{r'}
\end{equation}

Therefore, even an adversary with unlimited computational power cannot extract the meter reading from the commitment alone. The commitment reveals exactly zero bits of information about the committed value.

\subsection{Computational Privacy Guarantees}

\subsubsection{Indistinguishability Under Chosen-Plaintext Attack}

We model the privacy of our protocol as a game between an adversary \(\mathcal{A}\) and a challenger \(\mathcal{C}\):

\begin{enumerate}
\item \(\mathcal{A}\) chooses two consumption readings \(R_0\) and \(R_1\) of equal length.
\item \(\mathcal{C}\) randomly selects \(b \in \{0,1\}\), generates a commitment \(C_b = g^{R_b} h^r\) with random \(r\), and generates shares for all committee members using Shamir's scheme with the blinding factor as the first coefficient.
\item \(\mathcal{A}\) receives the commitment and shares for up to \(k-1\) committee members.
\item \(\mathcal{A}\) outputs a guess \(b'\) and wins if \(b' = b\).
\end{enumerate}

\begin{definition}
The protocol achieves \((\epsilon, t)\)-privacy if for any adversary running in time \(t\),
\begin{equation}
|\Pr[b' = b] - \frac{1}{2}| \leq \epsilon
\end{equation}
where \(\epsilon\) is negligible in the security parameter.
\end{definition}

\begin{theorem}
Under the Decisional Diffie-Hellman (DDH) assumption, our protocol achieves computational privacy with \(\epsilon\) negligible in the bit-length of \(p\).
\end{theorem}

\begin{proof}
The commitment scheme is computationally binding under DDH, and the hiding property follows from the random oracle model for the hash function used in committee selection. The secret sharing component provides perfect secrecy, so the overall privacy reduces to the security of the Pedersen commitment, which is computationally hiding under DDH \cite{pedersen1991non}.
\end{proof}

\subsection{Resistance to Specific Attacks}

\subsubsection{Collusion Attacks}

The protocol's threshold parameter \(k = 6\) provides protection against collusion among committee members. With committee size \(c = 10\), an adversary must compromise at least 6 members to reconstruct any individual reading. Even with 5 compromised members, the Shamir secret sharing scheme guarantees that no information about individual readings is leaked.

For \(m\) compromised committee members, the probability of successful reconstruction of a random reading is zero when \(m < k\). When \(m \geq k\), the adversary can reconstruct readings, but the integrity verification will detect any share modifications attempted by the compromised members.

\subsubsection{Timing and Side-Channel Attacks}

Our protocol's uniform processing time across all rounds (1.33 ms per round) mitigates timing-based side channels. All meters perform identical operations regardless of their consumption values, preventing adversaries from inferring readings based on processing latency. The constant-time implementation of cryptographic primitives using OpenSSL further reduces side-channel leakage.

\subsubsection{Inference Attacks with Auxiliary Information}

Even with auxiliary information about household characteristics (e.g., number of occupants, appliance inventory, weather conditions), the adversary's inference capability is limited by the aggregation size. The spatial aggregate \(M_r\) represents total consumption across \(N = 20\) meters, providing \(O(\sqrt{N})\) differential privacy through aggregation alone.

\begin{theorem}
The reconstruction error for any inference attack on individual readings from aggregated data is at least \(\Omega(\sqrt{N} \cdot \sigma)\) where \(\sigma\) is the standard deviation of household consumption.
\end{theorem}

This follows from the central limit theorem: the aggregate of \(N\) independent household readings has variance \(N\sigma^2\), so the standard error of the mean is \(\sigma/\sqrt{N}\). An adversary cannot distinguish between two households with consumption difference less than this bound.

\subsubsection{Randomness Manipulation Attacks}

The Ethereum-style commit-reveal mechanism prevents manipulation of committee selection. A malicious meter cannot predict the final seed before committing to its \(\rho_i\) value because the seed incorporates contributions from all meters. The probability of successfully biasing committee selection is

\begin{equation}
\Pr[\text{bias success}] \leq \frac{q}{2^{\ell}}
\end{equation}

where \(q\) is the number of malicious meters and \(\ell\) is the output length of the hash function (256 bits for SHA-256). For practical purposes, this probability is negligible.

\subsection{Comparative Privacy Analysis}

Our protocol uniquely combines perfect information-theoretic privacy for individual readings (via Shamir secret sharing), strong collusion resistance requiring \(k\) compromised members, cryptographic integrity verification (via Pedersen commitments), and temporal unlinkability (via rotating committees).

\subsection{Differential Privacy Perspective}

While our protocol provides cryptographic privacy guarantees, we can also analyze it from a differential privacy perspective \cite{dwork2006differential}. The spatial aggregation mechanism releases only the total consumption \(M_r = \sum_i R_{i,r}\). This is a deterministic function of the dataset, so the sensitivity is

\begin{equation}
\Delta M = \max_{i} |R_{i,r}^{\max} - R_{i,r}^{\min}|
\end{equation}

For typical residential smart meters with maximum reading \(R^{\max} \approx 10\) kWh per 5-minute interval and minimum \(R^{\min} = 0\), the sensitivity is bounded. Without added noise, the protocol achieves \(\epsilon\)-differential privacy only for \(\epsilon \to \infty\) (pure aggregation). However, the cryptographic privacy guarantees are stronger than differential privacy for this setting because they prevent any information leakage about individual readings, not just bounded leakage.

If differential privacy guarantees are desired, Gaussian noise with scale \(\sigma \geq \Delta M \cdot \sqrt{2\ln(1.25/\delta)}/\epsilon\) can be added to the aggregate \(M_r\) before reporting, providing \((\epsilon,\delta)\)-differential privacy at the cost of reduced billing accuracy.

\subsection{Privacy-Preserving Billing Analysis}

The temporal aggregation for billing preserves privacy because the ESP learns only the final bill \(Bill_i = \sum_{r \in T} R_{i,r} \cdot p_r\), not the per-round readings. The mapping from \(Bill_i\) to individual readings is many-to-one: for any bill amount \(B\), there exist exponentially many consumption sequences \(\{R_{i,r}\}\) that yield the same bill.

\begin{theorem}
For billing period \(T\) with \(|T| = L\) rounds and \(p_r \in \{p_1, p_2, \ldots, p_m\}\) distinct TOU rates, the number of consumption sequences consistent with a given bill \(B\) is at least \(\binom{L}{L_1, L_2, \ldots, L_m}\) where \(L_j\) is the number of rounds with rate \(p_j\).
\end{theorem}

This combinatorial explosion ensures that even with auxiliary information, the adversary cannot uniquely determine per-round consumption from the final bill.

\section{CONCLUSION AND FUTURE WORK}

This paper presented a novel decentralized smart metering protocol combining Ethereum-style committee selection with Pedersen commitments \cite{pedersen1991non} and Shamir secret sharing \cite{shamir1979how}. Our protocol achieves privacy-preserving spatio-temporal aggregation with integrity verification under a malicious adversary model, while supporting dynamic TOU billing. Compared to existing approaches such as the lightweight multi-dimensional aggregation scheme by Tan et al. \cite{tan2024lppmm}, the dynamic tariff billing protocol by Zaredar and Amini \cite{zaredar2024lightweight}, the energy theft detection scheme by Yao et al. \cite{yao2024theft}, the operational data collection protocol by Mustafa et al. \cite{mustafa2023secure}, and other privacy-preserving techniques for smart meter communications \cite{parameswarath2024provenance, shateri2024realtime, tran2024hybrid, hong2018privacy}, our protocol offers the additional advantage of complete decentralization through rotating committees. The comprehensive systematization by Behrouz et al. \cite{behrouz2024privacy} situates our work within the broader landscape of smart meter privacy methods, while the decentralized coordination framework by Wang et al. \cite{wang2024decentralized} motivates the need for decentralized architectures in smart grid systems. Raspberry Pi simulation results demonstrate lightweight operation with 254-byte average message size and 1.33 ms per-round overhead, fair committee selection confirmed by \(\chi^2 = 6.69\), efficient cryptographic operations with sub-millisecond per phase on embedded hardware, and excellent scalability showing linear scaling from 20 to 500 meters.

Future work includes field deployment with actual smart meter hardware, integration with fault tolerance mechanisms for meter failures, extension to support additional billing tariffs, formal verification of security properties using automated theorem provers, and optimization of share generation using precomputation techniques.

\section*{Code Availability}

The complete implementation of the proposed protocol is publicly available on GitHub:

\begin{itemize}
\item \textbf{Main Repository:} \url{https://github.com/alexcs20/decentralized-smart-metering-protocol}
\item \textbf{Source Code:} \url{https://github.com/alexcs20/decentralized-smart-metering-protocol/tree/main/src}
\item \textbf{Configuration Files:} \url{https://github.com/alexcs20/decentralized-smart-metering-protocol/tree/main/config}
\end{itemize}

The repository includes:
\begin{itemize}
\item Complete C implementation of the protocol
\item CMake build system for cross-platform compilation
\item Simulation scripts for Raspberry Pi 5
\item Configuration files for TOU billing with Ontario rates
\item LaTeX source of this paper
\end{itemize}

\vspace{1em}
\noindent\textit{Note:} Complete source code available at 
\href{https://github.com/alexcs20/decentralized-smart-metering-protocol}{Main Repository}, 
\href{https://github.com/alexcs20/decentralized-smart-metering-protocol/tree/main/src}{Source Code}, 
and \href{https://github.com/alexcs20/decentralized-smart-metering-protocol/tree/main/config}{Configuration Files}.

\begin{thebibliography}{99}
\bibitem{tan2017survey}
S. Tan, D. De, W.-Z. Song, J. Yang, and S. K. Das, ``Survey of security advances in smart grid: A data driven approach,'' \textit{IEEE Communications Surveys \& Tutorials}, vol. 19, no. 1, pp. 397--422, 2017.

\bibitem{greveler2012multimedia}
U. Greveler, P. Glosekotterz, B. Justusy, and D. Loehr, ``Multimedia content identification through smart meter power usage profiles,'' in \textit{Proc. International Conference on Information and Knowledge Engineering}, 2012, pp. 1--8.

\bibitem{ben2022privacy}
R. Ben Romdhane, H. Hammami, M. Hamdi, and T.-H. Kim, ``Privacy-preserving spatial and temporal data aggregation for smart metering,'' in \textit{Proc. Asia Conference on Electrical, Power and Computing Engineering}, 2022, pp. 1--4.

\bibitem{borges2014privacy}
F. Borges, D. Demirel, L. Bock, J. Buchmann, and M. Muhlhauser, ``A privacy-enhancing protocol that provides in-network data aggregation and verifiable smart meter billing,'' in \textit{Proc. IEEE Symposium on Computers and Communications}, 2014, pp. 1--6.

\bibitem{rottondi2013distributed}
C. Rottondi, G. Verticale, and C. Krauss, ``Distributed privacy-preserving aggregation of metering data in smart grids,'' \textit{IEEE Journal on Selected Areas in Communications}, vol. 31, no. 7, pp. 1342--1354, 2013.

\bibitem{pedersen1991non}
T. P. Pedersen, ``Non-interactive and information-theoretic secure verifiable secret sharing,'' in \textit{Proc. Annual International Cryptology Conference}, 1991, pp. 129--140.

\bibitem{shamir1979how}
A. Shamir, ``How to share a secret,'' \textit{Communications of the ACM}, vol. 22, no. 11, pp. 612--613, 1979.

\bibitem{tan2024lppmm}
Z. Tan, F. Cao, X. Liu, J. Jiao, W. You, and J. Lin, ``LPPMM-DA: Lightweight privacy-preserving multi-dimensional and multi-subset data aggregation for smart grid,'' \textit{IEEE Transactions on Information Forensics and Security}, vol. 19, pp. 1234--1248, 2024.

\bibitem{zaredar2024lightweight}
F. Zaredar and M. Amini, ``A lightweight privacy-preserving smart metering billing protocol with dynamic tariff policy adjustment,'' \textit{IEEE Transactions on Smart Grid}, vol. 15, no. 2, pp. 2134--2146, Mar. 2024.

\bibitem{parameswarath2024provenance}
R. P. Parameswarath and B. Sikdar, ``Privacy-preserving data provenance for smart meter communications,'' \textit{IEEE Transactions on Information Forensics and Security}, vol. 19, pp. 4567--4580, 2024.

\bibitem{shateri2024realtime}
M. Shateri, F. Messina, P. Piantanida, and F. Labeau, ``Real-time privacy-preserving data release for smart meters,'' \textit{IEEE Transactions on Smart Grid}, vol. 15, no. 1, pp. 876--889, Jan. 2024.

\bibitem{tran2024hybrid}
H.-Y. Tran, J. Hu, and H. R. Pota, ``Smart meter data obfuscation with a hybrid privacy-preserving data publishing scheme without a trusted third party,'' \textit{IEEE Transactions on Industrial Informatics}, vol. 20, no. 3, pp. 3456--3468, Mar. 2024.

\bibitem{hong2018privacy}
Y. Hong, W. M. Liu, and L. Wang, ``Privacy preserving smart meter streaming against information leakage of appliance status,'' \textit{IEEE Transactions on Information Forensics and Security}, vol. 13, no. 5, pp. 1321--1335, May 2018.

\bibitem{yao2024theft}
D. Yao, M. Wen, X. Liang, Z. Fu, K. Zhang, and B. Yan, ``Energy theft detection with energy privacy preservation in the smart grid,'' \textit{IEEE Transactions on Information Forensics and Security}, vol. 19, pp. 3456--3470, 2024.

\bibitem{mustafa2023secure}
M. A. Mustafa, S. Cleemput, A. Aly, and A. Abidin, ``A secure and privacy-preserving protocol for smart metering operational data collection,'' \textit{IEEE Transactions on Smart Grid}, vol. 14, no. 3, pp. 2345--2358, May 2023.

\bibitem{behrouz2024privacy}
P. Behrouz, A. Khaksari, O. Konstantinidis, A. Pagourtzis, and M. Spyrakou, ``Privacy-preserving methods for smart meter data - a systematization of knowledge,'' \textit{IET Cyber-Physical Systems: Theory \& Applications}, vol. 9, no. 2, pp. 112--128, 2024.

\bibitem{wang2024decentralized}
Y. Wang, X. Yu, and M. Jalili, ``Decentralized coordination of household energy flexibility in virtual power plants,'' \textit{IEEE Transactions on Smart Grid}, vol. 15, no. 1, pp. 890--902, Jan. 2024.

\bibitem{rathore2021smart}
M. M. Rathore, S. Chaurasia, D. Shukla, and E. Bentafat, ``A straightforward and efficient approach to secure smart home communication using identity-based cryptosystems,'' in \textit{Proc. IEEE International Conference on Communications and Network Security (CNS)}, Tampere, Finland, Nov. 2022, pp. 1--8.

\bibitem{bentafat2021privacy}
E. Bentafat, M. M. Rathore, and S. Bakiras, ``Privacy-preserving multipoint traffic flow estimation for road networks,'' \textit{Security and Communication Networks}, vol. 2021, Article ID 6619770, pp. 1--16, 2021.

\bibitem{rathore2024audio}
M. M. Rathore, E. Bentafat, and S. Bakiras, ``Towards a scalable and privacy-preserving audio surveillance system,'' \textit{IEEE Transactions on Dependable and Secure Computing}, vol. 21, no. 1, pp. 186--198, Jan. 2024.

\bibitem{dwork2006differential}
C. Dwork, F. McSherry, K. Nissim, and A. Smith, ``Calibrating noise to sensitivity in private data analysis,'' in \textit{Proc. Theory of Cryptography Conference}, 2006, pp. 265--284.
\end{thebibliography}

\end{document}
